\clearpage\pagestyle{otMain}\fontsize{11}{15}\selectfont
\corrigesChapitre{6}{Les fonctions de référence}
\begin{blocLivre}{cor-6-1}
\begin{corrige}{1}{Une fonction affine}
\demarche{Lire les coefficients puis utiliser deux points distincts.}
La pente est $\res{-2}$, l’ordonnée à l’origine est $\res3$. On a $f(4)=-8+3=\res{-5}$. Pour tracer, utiliser par exemple $(0;3)$ et $(1;1)$. Lorsque l’abscisse augmente de 1, l’ordonnée diminue de 2, conformément à la pente.
\label{fig-cor-6-1}
\figureGeo[corrige]{reference-affine}
\end{corrige}
\end{blocLivre}
\begin{blocLivre}{cor-6-2}
\begin{corrige}{2}{Comparer des carrés}
\demarche{Localiser les entrées sur un intervalle de monotonie.}
Les distances à 0 valent 7 et 3, donc $\res{(-7)^2>(-3)^2}$. Sur les nombres positifs, la fonction carré est croissante : $0{,}4<0{,}7$ entraîne $\res{0{,}4^2<0{,}7^2}$. Les variations dépendent du côté de 0 où se trouvent les nombres.
\end{corrige}
\end{blocLivre}
\begin{blocLivre}{cor-6-3}
\begin{corrige}{3}{La fonction inverse}
\demarche{Distinguer absence d’image et absence d’antécédent.}
Les images sont respectivement $\res{-1/4}$, $\res2$ et $\res{1/5}$. Le nombre 0 n’a aucun antécédent, car $1/x=0$ impliquerait $1=0$ après multiplication par un $x$ non nul. Attention : 0 n’a pas non plus d’image, puisque la fonction n’y est pas définie.
\label{fig-cor-6-3}
\figureGeo[corrige]{reference-inverse}
\end{corrige}
\end{blocLivre}
\begin{blocLivre}{cor-6-4}
\begin{corrige}{4}{Racine et domaine}
\demarche{Examiner les signes avant de mettre au carré.}
Le domaine est $x\ge0$. Dans la première équation, les deux membres sont positifs : on peut élever au carré, ce qui donne $\res{x=25}$. Dans la seconde, le membre de gauche est positif ou nul et celui de droite négatif : il n’y a aucune solution. Élever au carré trop tôt aurait créé la fausse solution 4.
\label{fig-cor-6-4}
\figureGeo[corrige]{reference-racine}
\end{corrige}
\end{blocLivre}
\begin{blocLivre}{cor-6-5}
\begin{corrige}{5}{La courbe en V}
\demarche{Interpréter la valeur absolue comme distance à zéro.}
Les trois valeurs sont $\res3$, $\res0$ et $\res{2{,}5}$. L’inégalité demande les nombres à distance au plus 4 de 0 : $\res{x\in[-4;4]}$. On retrouve sur la courbe en V les abscisses pour lesquelles l’ordonnée ne dépasse pas 4.
\label{fig-cor-6-5}
\figureGeo[corrige]{reference-absolue}
\end{corrige}
\end{blocLivre}
\begin{blocLivre}{cor-6-6}
\begin{corrige}{6}{Une équation de cube}
\demarche{Utiliser la croissance de la fonction cube sur toute la droite réelle.}
Comme $(-3)^3=-27$ et que la fonction cube est strictement croissante, l’unique solution est $\res{x=-3}$. La même croissance donne $-1{,}2<-1{,}1\Rightarrow\res{(-1{,}2)^3<(-1{,}1)^3}$. Il n’y a pas de changement de sens pour le cube sur les négatifs.
\label{fig-cor-6-6}
\figureGeo[corrige]{reference-cube}
\end{corrige}
\end{blocLivre}
\begin{blocLivre}{cor-6-7}
\begin{corrige}{7}{Une inéquation de carré}
\demarche{Lire les deux branches de la parabole.}
Les valeurs frontières sont $-3$ et 3. Le carré est inférieur à 9 entre ces bornes : $\res{]-3;3[}$. Il est supérieur ou égal à 9 à l’extérieur : $\res{]-\infty;-3]\cup[3;+\infty[}$. Les solutions négatives sont aussi nécessaires que les positives.
\label{fig-cor-6-7}
\figureGeo[corrige]{reference-carre}
\end{corrige}
\end{blocLivre}
\begin{blocLivre}{cor-6-8}
\begin{corrige}{8}{Comparer des inverses}
\demarche{Utiliser la monotonie sur un intervalle où la fonction est définie.}
Sur $]0;+\infty[$, l’inverse est décroissante : $3<7$ donne $\res{1/3>1/7}$. Sur $]-\infty;0[$, elle l’est aussi : $-5<-2$ donne $\res{-1/5>-1/2}$. On traite les paires séparément car 0 coupe le domaine ; on ne peut pas annoncer une décroissance sur toute la réunion des deux intervalles.
\end{corrige}
\end{blocLivre}
\begin{blocLivre}{cor-6-9}
\begin{corrige}{9}{Carré et identité}
\demarche{Comparer par différence, puis exploiter la contrainte x positif.}
On factorise $x^2-x=x(x-1)$. Sur $x\ge0$, ce produit est négatif ou nul exactement pour $0\le x\le1$. La réponse est $\res{[0;1]}$. Les deux égalités sont obtenues en $\res{x=0\text{ et }x=1}$. Le dessin illustre les positions relatives des deux courbes.
\label{fig-cor-6-9}
\figureGeo[corrige]{carre-identite}
\end{corrige}
\end{blocLivre}
\begin{blocLivre}{cor-6-10}
\begin{corrige}{10}{Une inéquation d’inverse}
\demarche{Ramener le quotient à zéro et étudier les signes.}
On impose $x\ne0$ puis on écrit $(1-2x)/x\ge0$. Les valeurs à examiner sont 0, interdite, et $1/2$, zéro du numérateur. Le quotient est positif sur $]0;1/2[$ et nul en $1/2$. Ainsi $\res{S=]0;1/2]}$. Les négatifs ne conviennent pas car leur inverse est négatif.
\end{corrige}
\end{blocLivre}
\begin{blocLivre}{cor-6-11}
\begin{corrige}{11}{Une distance minimale}
\demarche{Lire la formule comme une distance, puis comparer les extrémités.}
$f(x)$ mesure la distance de $x$ à 3. Elle est minimale en 3, avec $\res{f(3)=0}$. Les distances aux extrémités sont 3 et 5 ; le point le plus éloigné est 8, donc le maximum est $\res5$, atteint en $\res{x=8}$. Sur cet intervalle, la fonction décroît jusqu’à 3 puis croît.
\end{corrige}
\end{blocLivre}
\begin{blocLivre}{cor-6-12}
\begin{corrige}{12}{Un problème de surface}
\demarche{Passer d’une aire à une longueur par la racine, pas par une division.}
Notons $c\ge0$ le côté. $25\le c^2\le81$ donne $\res{c\in[5;9]}$ cm par croissance de la racine. Le périmètre vaut $4c$, donc $\res{P\in[20;36]}$ cm. Les deux bornes sont atteignables avec des carrés d’aires 25 et 81 cm$^2$.
\end{corrige}
\end{blocLivre}
\begin{blocLivre}{cor-6-13}
\begin{corrige}{13}{Une équation mixte}
\demarche{Fixer le domaine et les signes avant de comparer les carrés.}
Le domaine impose $x\ge0$ ; les deux membres sont alors positifs ou nuls. L’égalité équivaut à $x=x^2$, soit $x(x-1)=0$. On obtient $\res{S=\{0;1\}}$. Les deux valeurs respectent le domaine et satisfont l’équation initiale. La condition de signe rendait l’élévation au carré réversible.
\end{corrige}
\end{blocLivre}
\begin{blocLivre}{cor-6-14}
\begin{corrige}{14}{Une affirmation sur les fonctions}
\demarche{Vérifier les hypothèses de monotonie au lieu de retenir une règle sans domaine.}
Pour les carrés, $-3<-2$ mais $9>4$ : la phrase est fausse sur tous les réels ; elle est vraie pour $0\le a<b$. Pour les inverses, $-1<1$ mais $-1<1$ aussi après inversion : le sens annoncé échoue lorsque l’on traverse 0. Il est correct si $a,b$ appartiennent tous deux à $]0;+\infty[$ ou tous deux à $]-\infty;0[$.
\end{corrige}
\end{blocLivre}
\begin{blocLivre}{cor-6-15}
\begin{corrige}{15}{Un encadrement transformé}
\demarche{Appliquer la monotonie et vérifier l’atteinte des bornes.}
Sur les positifs, carré et racine sont croissantes : $\res{4\le x^2\le25}$ et $\res{\sqrt2\le\sqrt x\le\sqrt5}$. L’inverse est décroissante : $\res{1/5\le1/x\le1/2}$. Chaque borne est atteinte pour $x=2$ ou $x=5$ ; ces encadrements sont donc les plus précis possibles à partir de l’information donnée.
\end{corrige}
\end{blocLivre}
